\documentclass[../main.tex]{subfiles}

\begin{document}

\subsection{Overview}

CVE-2024-42072 is a \textbf{Denial-of-Service vulnerability} affecting the Linux eBPF subsystem in kernel versions 6.9 up to 6.9.8. The bug is related to the newly introduced \textbf{\texttt{may\_goto} instruction}, which was added to support bounded loops in eBPF programs.

Historically, the eBPF verifier rejected unbounded loops to guarantee that every BPF program terminates. The \texttt{may\_goto} instruction was introduced as a controlled mechanism for loops: it relies on a hidden counter stored on the BPF stack. At each iteration, the counter is decremented. When the counter reaches zero, execution should exit the loop.

The vulnerability appears when \texttt{may\_goto} is used with a \textbf{negative offset}, i.e., a backward jump. Due to an \textbf{incorrect patching formula} and \textbf{insufficient verifier loop-pruning checks}, a BPF program containing only a few instructions can be accepted by the verifier but transformed at runtime into an \textbf{infinite kernel loop}.

\begin{table}[H]
\centering
\begin{tabularx}{\textwidth}{|l|X|}
\hline
\textbf{CVE ID} & CVE-2024-42072 \\
\hline
\textbf{Affected versions} & Linux kernel 6.9--6.9.8 \\
\hline
\textbf{Vulnerability type} & eBPF verifier / JIT fixup logic error \\
\hline
\textbf{Root cause} & Incorrect \texttt{may\_goto} patching for negative offsets and unsafe state pruning \\
\hline
\textbf{Impact} & Denial of Service, kernel soft lockup \\
\hline
\textbf{CVSS} & 7.8 HIGH \\
\hline
\end{tabularx}
\caption{Summary of CVE-2024-42072.}
\end{table}

Unlike the other vulnerabilities analyzed in this report, this CVE does not provide a memory disclosure or arbitrary write primitive. Its impact is purely availability-oriented: once triggered, the target CPU core remains trapped inside a kernel loop and the system becomes unresponsive, requiring a hard reboot.

\subsection{Technical Background: \texttt{may\_goto}}

The \texttt{may\_goto} instruction is not directly executed as a native BPF instruction. Before execution, the kernel expands it into a short sequence of ordinary BPF instructions implementing a bounded-loop counter.

Conceptually, \texttt{may\_goto} behaves as follows:

\begin{lstlisting}[caption={Expected \texttt{may\_goto} semantics.}]
if counter > 0:
    counter -= 1
    jump to target
else:
    continue after the loop
\end{lstlisting}

Internally, the kernel rewrites it into a sequence similar to:

\begin{lstlisting}[caption={Simplified \texttt{may\_goto} expansion.}]
AX = *(fp + counter_offset)
if AX == 0:
    jump to loop_exit
AX -= 1
*(fp + counter_offset) = AX
\end{lstlisting}

The security of this mechanism depends on two properties:

\begin{enumerate}
    \item the patching phase must compute the \textbf{correct jump target};
    \item the verifier must still \textbf{detect infinite-loop patterns}.
\end{enumerate}

CVE-2024-42072 breaks both properties.

\subsection{Vulnerable Code Logic}

The vulnerability is caused by the combination of two distinct bugs.

\subsubsection{Bug 1: Incorrect Negative Offset Patching}

The first bug is located in \texttt{do\_misc\_fixups()}, the function responsible for expanding \texttt{may\_goto} into ordinary BPF instructions. The vulnerable code used the same offset formula for both forward and backward jumps:

\begin{lstlisting}[language=C, caption={Vulnerable offset calculation.}]
insn_buf[1] = BPF_JMP_IMM(BPF_JEQ,
                          BPF_REG_AX,
                          0,
                          insn->off + 2);
\end{lstlisting}

This formula is valid for positive offsets, but wrong for negative ones. For example, with:

\begin{lstlisting}
may_goto -4
\end{lstlisting}

the expected behavior is that, when the counter reaches zero, execution jumps to the loop exit block. Instead, the vulnerable formula computes:

\begin{lstlisting}
JEQ offset = -4 + 2 = -2
\end{lstlisting}

This makes the generated conditional jump target the counter-loading instruction itself. When the counter becomes zero, execution loops forever between:

\begin{lstlisting}
AX = *(fp + counter_offset)
if AX == 0: jump back to AX load
\end{lstlisting}

Therefore, the bounded loop degenerates into an \textbf{infinite loop}.

\subsubsection{Bug 2: Unsafe State Pruning}

The second bug is in the verifier state-pruning logic, specifically around \texttt{is\_state\_visited()}. The verifier avoids path explosion by skipping states already considered equivalent. For \texttt{may\_goto}, the vulnerable logic could prune a state without properly accounting for the current \texttt{may\_goto\_depth}.

The vulnerable behavior can be summarized as:

\begin{lstlisting}[caption={Simplified vulnerable pruning logic.}]
if instruction is may_goto:
    if current_state equals old_state:
        prune path
    skip infinite-loop check
\end{lstlisting}

This prevents the verifier from detecting that the accepted program may not terminate after the buggy patching transformation. The result is a verifier-approved program that becomes non-terminating only after \textbf{runtime fixups}.

\subsection{Exploit Technical Background}

The exploit uses a minimal BPF program of seven instructions. Its purpose is not to corrupt memory, but to force the vulnerable kernel to execute the \textbf{wrongly patched \texttt{may\_goto} path}.

The logical program is:

\begin{lstlisting}[caption={Minimal exploit structure.}]
0: r0 = 0
1: goto +1
2: exit
3: r1 = 0
4: r1 += 1
5: may_goto -4
6: goto -3
\end{lstlisting}

Instruction 5 is the critical one. It should behave as a bounded backward loop. Once the hidden counter reaches zero, execution should jump to instruction 2 and terminate.

However, on a vulnerable kernel, \texttt{may\_goto -4} is expanded with the wrong conditional-jump target:

\begin{lstlisting}[caption={Buggy runtime expansion.}]
5: AX = *(fp + counter_offset)
6: if AX == 0 goto -2
7: AX -= 1
8: *(fp + counter_offset) = AX
9: goto -6
\end{lstlisting}

The unconditional jump at instruction 9 corresponds to the original \texttt{goto -3}, adjusted after the \texttt{may\_goto} expansion inserted three additional instructions. Its purpose is still to jump back to the loop body.

The real bug is in instruction 6. The conditional jump should leave the loop when the counter reaches zero. Instead, due to the \textbf{incorrect offset calculation}, its target becomes instruction 5. Therefore, once \texttt{AX == 0}, the program repeatedly executes:

\begin{lstlisting}
5: AX = *(fp + counter_offset)
6: if AX == 0 goto 5
\end{lstlisting}

This creates a \textbf{deterministic infinite loop} inside kernel context.

\subsection{Exploit Execution}

The userspace exploit performs only the minimal operations required to load and trigger the BPF program:

\begin{enumerate}
    \item load the crafted BPF program with \texttt{bpf(BPF\_PROG\_LOAD)};
    \item create a socket pair;
    \item attach the BPF program with \texttt{setsockopt(SO\_ATTACH\_BPF)};
    \item send data through the socket with \texttt{write()}.
\end{enumerate}

The write operation triggers execution of the socket-filter BPF program. Once the vulnerable \texttt{may\_goto} path is reached, the CPU core remains trapped in the infinite loop.

\begin{figure}[H]
    \centering
    \includegraphics[width=\textwidth]{img/CVE-2024-42072.png}
    \caption{Exploit execution log for CVE-2024-42072.}
    \label{fig:log-2024-42072}
\end{figure}

After the trigger, the virtual machine becomes unresponsive. In the QEMU environment, recovery requires killing the VM process from the host. On physical hardware, the same condition would require a hard reboot.

\subsection{Impact Analysis}

The vulnerability does not expose confidentiality or integrity primitives. It does not leak kernel pointers, does not provide arbitrary read/write, and does not directly support privilege escalation.

Its impact is instead a pure availability failure:

\begin{itemize}
    \item the verifier accepts a program that should be rejected;
    \item runtime fixups transform \texttt{may\_goto} into an infinite loop;
    \item the loop executes in kernel context;
    \item the affected CPU core stops making progress;
    \item the system becomes unresponsive.
\end{itemize}

This makes CVE-2024-42072 different from memory-corruption eBPF exploits. The attacker does not need heap grooming, KASLR bypass, or post-exploitation primitives. The exploit is extremely small and reliable: the security failure is entirely in control-flow validation and instruction rewriting.

\subsection{Patch and Mitigation}

The fix is associated with the upstream/stable patch series referenced for this CVE, including:

\begin{lstlisting}
175827e04f4be53f3dfb57edf12d0d49b18fd939
2b2efe1937ca9f8815884bd4dcd5b32733025103
\end{lstlisting}

The patch addresses both vulnerable components.

First, the patching logic now treats negative \texttt{may\_goto} offsets differently from positive ones. This is necessary because the exit jump generated during \texttt{may\_goto} expansion must target a different instruction depending on whether the original jump is forward or backward:

\begin{lstlisting}[language=C, caption={Conceptual corrected offset calculation.}]
if (insn->off >= 0)
    exit_off = insn->off + 2;
else
    exit_off = insn->off - 1;

insn_buf[1] = BPF_JMP_IMM(BPF_JEQ,
                          BPF_REG_AX,
                          0,
                          exit_off);
\end{lstlisting}

For \texttt{may\_goto -4}, the corrected formula produces:

\begin{lstlisting}
JEQ offset = -4 - 1 = -5
\end{lstlisting}

which points to the intended exit block instead of looping back to the counter load.

Second, the verifier pruning logic was changed so that a state is not pruned when the current \texttt{may\_goto\_depth} is equal to the depth stored in the previously visited state. In that case, the verifier has reached the same \texttt{may\_goto} state again, which indicates an actual infinite loop rather than a safe repeated state. The verifier must therefore continue exploration until it reaches \texttt{bpf\_exit}.

\begin{lstlisting}[caption={Conceptual model of the pruning fix.}]
if states_equal:
    if cur.may_goto_depth == visited.may_goto_depth:
        reject_or_continue_infinite_loop_check
    else:
        prune safely
else:
    continue verification
\end{lstlisting}

\subsection{Security Impact}

CVE-2024-42072 highlights a different class of eBPF verifier failures. Many verifier vulnerabilities lead to memory corruption by producing incorrect bounds or stale pointer states. Here, the failure is about guaranteed termination.

The verifier's role is not only to prevent illegal memory accesses, but also to ensure that BPF programs cannot monopolize kernel execution. The \texttt{may\_goto} instruction was designed to relax the historical ban on loops while preserving that guarantee. CVE-2024-42072 shows that introducing new control-flow features into eBPF is dangerous unless every transformation performed after verification is semantically equivalent to what the verifier analyzed.

In short, the verifier accepted a bounded loop, but the kernel executed an unbounded one.

\end{document}